%
% Cover letter for AIP Advances submission
% (see AIP_Style_4thed.pdf, Sec. II.D, on including a cover letter)
%
\documentclass[11pt]{letter}

\usepackage[T1]{fontenc}
\usepackage{lmodern}
\usepackage[margin=1in,top=0.75in,bottom=0.75in]{geometry}

% Sender address block (corresponding author)
\address{
Hideo Takahashi \\
Managed \& Platform Services Business Division, Hitachi, Ltd. \\
292, Yoshida, Totsuka-ku, Yokohama, 244-0817, Japan \\
hideo.takahashi.qa@hitachi.com
}

\signature{Hideo Takahashi}

\begin{document}

\begin{letter}{The Editors \\ AIP Advances \\ AIP Publishing}

\opening{Dear Editors,}

We are pleased to submit our manuscript, entitled ``Grid-Based Quantum Circuit
Simulation of the Hydrogen Atom: Coulomb Singularity Regularization and
Parameter-Selection Guidelines,'' by Hideo Takahashi, Tatsuya Tomaru,
Toshiyuki Hirano, Saisei Tahara, and Fumitoshi Sato, for consideration as a
Research Article in AIP Advances.

Quantum computers are expected to revolutionize quantum chemistry simulation.
Among the various methods, the grid-based first-quantization approach is
well suited to chemical reactions for its ability to represent structural
changes in the wave function.

In this work, we address three challenges in grid-based first-quantization
simulation: (1) circuit design of the Coulomb potential under qubit constraints, (2)
mitigation of the Coulomb potential's singularity, and (3) selection of
discretization parameters that minimize error. (1) Uniformly Controlled
Rotation Z (UCRZ) circuits let us construct a Coulomb-potential circuit
compact enough for a small-scale emulator; although this approach does not
scale to larger models, it achieves, to our knowledge, the first complete
grid-based simulation of the Coulomb potential run end-to-end on a quantum
computer emulator, for simplified models such as the 1D/2D hydrogen atom.
(2) We regularize
the singularity via a geometric average of the potential near the pole, which
numerically minimizes the resulting energy error. (3) We determine the grid
spacing from the trade-off between spatial discretization and domain-truncation
error, and the time step from a sampling-theorem bound on the wave function's
frequency content rather than the worst-case Trotter error bound.

While (1) is specific to this small-scale, qubit-constrained regime, the
methods in (2) and (3) address the grid discretization itself and apply
broadly to grid-based first-quantization simulations at any scale. We
believe this combination of broadly applicable methodological guidance with
a concrete, emulator-verified demonstration will interest AIP Advances'
readership across quantum computing and computational physics.

Thank you for your consideration.

\closing{Sincerely,}

\end{letter}
\end{document}
